\labsection{4}{Factors, matrix operations, and multidimensional arrays}{sec:lab04}

\begin{labproject}{Lab 4b guided practice}
\textbf{Coverage.} Chapter~\ref{chap:factors-matrices-arrays}.\\
\textbf{Goal.} Summarize factor levels, compare a matrix with its transpose, and construct multidimensional arrays.

\textbf{Task 1 --- Retirement-age factor levels.}

\begin{lstlisting}[style=dsr]
age <- c(
  55,57,56,52,51,59,58,53,59,55,
  60,60,60,60,52,55,56,51,60,52,
  54,56,52,57,54,56,58,53,53,50,
  55,51,57,60,57,55,51,50,57,58
)

age_factor <- factor(age)
print(levels(age_factor))
print(table(age_factor))
\end{lstlisting}

The exact-age counts are:

\begin{dstable}
\begin{tabular}{@{}rrrrrrrrrrrr@{}}
\toprule
Age & 50 & 51 & 52 & 53 & 54 & 55 & 56 & 57 & 58 & 59 & 60 \\
\midrule
Count & 2 & 4 & 4 & 3 & 2 & 5 & 4 & 5 & 3 & 2 & 6 \\
\bottomrule
\end{tabular}
\end{dstable}

The requested ranges overlap at 52, 54, 56, and 58. Direct inclusive counts are:

\begin{lstlisting}[style=dsr]
range_counts <- c(
  `50-52` = sum(age >= 50 & age <= 52),
  `52-54` = sum(age >= 52 & age <= 54),
  `54-56` = sum(age >= 54 & age <= 56),
  `56-58` = sum(age >= 56 & age <= 58),
  `58-60` = sum(age >= 58 & age <= 60)
)
print(range_counts)
\end{lstlisting}

The counts are 10, 9, 11, 12, and 11. Do not total them because the groups overlap.

\begin{pitfall}
Non-overlapping groups require an explicit boundary convention; the source does not provide one.
\end{pitfall}

\textbf{Task 2 --- Create a 3-by-3 matrix and transpose it.}

\begin{lstlisting}[style=dsr]
V1 <- c(2, 3, 1, 5, 4, 6, 8, 7, 9)

Matrix1 <- matrix(V1, nrow = 3, byrow = TRUE)
Matrix2 <- t(Matrix1)

rownames(Matrix1) <- c("R1", "R2", "R3")
colnames(Matrix1) <- c("C1", "C2", "C3")
rownames(Matrix2) <- c("R1", "R2", "R3")
colnames(Matrix2) <- c("C1", "C2", "C3")

print(Matrix1 + Matrix2)
print(Matrix1 - Matrix2)
print(Matrix1 * Matrix2)
print(Matrix1 / Matrix2)
\end{lstlisting}

These operators are element-wise.

\textbf{Task 3 --- Construct the two arrays shown in the expected output.}

\begin{lstlisting}[style=dsr]
Array1 <- array(1:24, dim = c(2, 4, 3))
Array2 <- array(25:54, dim = c(3, 2, 5))

print(Array1)
print(Array2)

print("The second row of the second matrix of the first array:")
print(Array1[2, , 2])

print("The element corresponding to sample output 30:")
print(Array2[3, 2, 1])
\end{lstlisting}

\begin{pitfall}
Array2 has only two columns, although the task requests column 3. The expected value 30 corresponds to \texttt{Array2[3,2,1]}.
\end{pitfall}
\end{labproject}

\begin{labdeliverable}
Submit the factor summary, range interpretation, matrix results, arrays, and requested selections.
\end{labdeliverable}
